% ---------------------------------------------------------------------
\section{Ek bulgular}\label{sec:results-supp}
% ---------------------------------------------------------------------
% Tablolar ve şekiller \analysisdir (varsayılan ../../analysis) altından ya da
% tables/ ve figures/ içindeki yerel kopyalardan okunur. Metindeki her sayı
% kaynak hücresiyle birlikte paper/results_number_sources.md dosyasında
% listelenmiştir.

\TODO{Yazarlar raporlanan her değeri analiz çıktılarıyla karşılaştırarak denetler ve bu bölümün ifadelerini gözden geçirir}.

Bu bölümdeki tablolar analiz iş akışı tarafından üretilmektedir.
İş akışının T2, T3, T3b, T3c, T3d, T5, T4, T7, T6, TS1, TS2 ve TS3 dosyaları bu sırayla Tablo~\ref{tab:t2}--\ref{tab:ts3} olarak yer almaktadır.
Değerler tohum ortalamalı politikaların ortalamalarıdır; parantez içinde iki eğitim tohumu üzerindeki aralık verilmiştir ve \dag{} tek bir çalıştırmaya sahip deterministik yöntemleri işaretler.
Örnek gün şekli, planlanan ablasyonlar ve stres anma düzeyi dışındaki duyarlılık çalışmaları çalıştırılmamıştır.

\subsection{Öğrenme eğrileri ve sıfır-atış aktarım}\label{app:res-tables}

\Fig{fig:curves}, öğrenen yöntemlerin eğitim eğrilerini iki tohum üzerindeki bantlarla birlikte göstermektedir.
\Tab{tab:t2}, her istasyon boyutu için \Fig{M-fig:zeroshot} şeklinin dayandığı sıfır-atış sonuçlarını, $N=50$ için eğitilen referansa göre aktarım açığı dâhil olmak üzere listelemektedir.
Bu referans yalnızca $N=50$ için dağılım içidir; dolayısıyla $N=100$ ve $N=200$ için aktarım açığı iki dışdeğerlemeyi karşılaştırır.

\begin{figure}[t]
\centering
\resultfig{figS1_training_curves.pdf}{Öğrenen yöntemlerin eğitim eğrileri.}
\caption{Öğrenen yöntemlerin ortam adımlarına göre eğitim eğrileri; bantlar iki eğitim tohumu üzerindeki aralığı gösterir.}
\label{fig:curves}
\end{figure}

\restable{T2}{\placeholdertable{tab:t2}{İstasyon boyutları $N$ arasında sıfır-atış aktarım}{T2.tex}}

\subsection{Kısıtların sağlanması}\label{app:res-constraints}

Test ızgarası üzerinden birleştirilmiş olarak \Tab{tab:t3}; E4'ü, karşılanmayan enerjiyi, aşırı yükü, \Eq{eq:deficit} ile verilen açık sınırına göre eksik enerjiyi ve sınırlandırılmış ve kabul kontrolüyle düşürülmüş talepleri vermektedir.
Açık sınırı oranı yalnızca tam güvenlik katmanına sahip yöntemleri kapsar, çünkü ayrılış sınırı içermeyen varyantın böyle bir sınırı yoktur.
İstasyon boyutu başına \Tab{tab:t3b}; E4'ü, yedek kural durumlarına göre olursuz adımları ve toplam limitlerin ve ayrılış alt sınırının bağlayıcılık oranlarını vermektedir.
\Tab{tab:t3c} ve \Tab{tab:t3d}, her iki tohumun tüm epizotları üzerinden toplanan olay sayılarını genel olarak ve istasyon boyutu başına vermektedir.
Bu tabloların son sütunu, Teorem~\ref{M-thm:invariance}(b) için epizot düzeyindeki denetimdir: hiç olursuz adım içermeyen epizotlarda gereksinimi karşılanmayan EV'ler.
Bu sütun, ayrılış sınırına sahip konformal olmayan her yöntem için sıfırdır; ayrılış sınırı içermeyen varyantta ise bu tür 1.551 EV bulunmuştur.
Ancak tüm epizotlar üzerinden bu varyant, gereksinimi karşılanmayan daha az EV bırakmıştır: varyant için 1.551, \method{} için 6.158 (\Sec{M-sec:res-bound}).
Konformal varyantlar için denetim, yeniden sertifikalanmış gereksinimi ortak gereksinimin altında kalan ya da öngörülen sınırdan önce ayrılan EV'leri de saymaktadır; teorem bu durumları kapsamaz.

\restable{T3}{\placeholdertable{tab:t3}{Tüm test ızgarası yapılandırmaları üzerinden yöntem başına kısıtların sağlanması}{T3.tex}}
\restable{T3b}{\placeholdertable{tab:t3b}{İstasyon boyutu $N$ başına E4, yedek kural durumlarına göre olursuz adımlar ve bağlayıcılık oranları}{T3b.tex}}
\restable{T3c}{\placeholdertable{tab:t3c}{Değerlendirilen tüm epizotlar üzerinden yöntem başına olay sayıları}{T3c.tex}}
\restable{T3d}{\placeholdertable{tab:t3d}{Yöntem ve istasyon boyutu $N$ başına olay sayıları}{T3d.tex}}

\subsection{Garantinin maliyeti ve kısıt yönetimi}\label{app:res-ablation}

\Tab{tab:t5}, \Fig{M-fig:cost} şeklinin dayandığı garantinin maliyetini dikkat tabanlı ve MLP çiftleri için $J^{\sm{pen}}$ ve kâr cinsinden vermektedir.
\Fig{fig:ablation} ve \Tab{tab:t4}, mimarileri ve kısıt yönetimi varyantlarını eğitim boyutlarında ve sıfır-atış aktarımda karşılaştırmaktadır.
$N\ge50$ için \method{} yönteminin açığı 0,901, Deep Sets+safety layer'ın 0,674 ve Attention+safety layer without departure bound'un 0,770 olmuştur.
Katmanı kullanım aşamasında uygulayan ceza ile eğitilmiş politika 0,910'a, Attention+penalty 0,996'ya ve Attention+PPO-Lagrangian 1,138'e ulaşmıştır.

\restable{T5}{\placeholdertable{tab:t5}{İstasyon boyutu $N$ başına garantinin maliyeti}{T5.tex}}

\begin{figure}[t]
\centering
\resultfig{figS2_ablation.pdf}{Mimari ve kısıt yönetimi varyantlarının açığı.}
\caption{\Tab{tab:t4} içindeki mimari ve kısıt yönetimi varyantlarının eğitim boyutlarında ve sıfır-atış aktarımdaki açığı.}
\label{fig:ablation}
\end{figure}

\restable{T4}{\placeholdertable{tab:t4}{Mimari ve kısıt yönetimi}{T4.tex}}

\subsection{Konformal ayrılış sınırı}\label{app:res-conformal}

\Fig{fig:conformal} ve \Tab{tab:t7}, \method{}-CP yönteminin kapsamasını, ortak gereksinime göre eksik enerjisini, açığını ve şarj edilen enerjisini kapsama dışı kalma düzeyi $\alpha$'ya göre göstermektedir.
0,99, 0,95, 0,90 ve 0,80 hedefleri için elde edilen kapsama 0,992, 0,950, 0,902 ve 0,799 olmuştur.
Tüm epizotlar üzerinden toplandığında eksik enerji $\alpha=0.01$ için 771.853 kWh, $\alpha=0.05$ için 12.507 kWh ve bilinen ayrılış adımlarıyla 2.568 kWh olmuştur.
Karşılaştırma için kalibre edilmemiş nokta tahmini 0,899 kapsamaya ve 1,322 açığına ulaşmıştır.

\begin{figure}[t]
\centering
\resultfig{figS3_conformal_alpha.pdf}{$\alpha$'ya göre kapsama, eksik enerji, açık ve şarj edilen enerji.}
\caption{Konformal ayrılış sınırının ödünleşimi: kapsama dışı kalma düzeyi başına hedefine göre kapsama, ortak gereksinime göre eksik enerji, açık ve şarj edilen enerji.}
\label{fig:conformal}
\end{figure}

\restable{T7}{\placeholdertable{tab:t7}{Konformal ayrılış sınırı}{T7.tex}}

\subsection{Hesaplama süresi}\label{app:res-time}

\Fig{fig:inference} ve \Tab{tab:t6}, kontrol adımı başına çıkarım süresini ve çalıştırma başına eğitim süresini raporlamaktadır.
\method{} yönteminin bir eğitim çalıştırması 1,15 saat sürmüş; bir kontrol adımı $N=10$ için 1,68 ms, $N=200$ için 7,76 ms almıştır.
$N=200$ için Deep Sets+safety layer adım başına 5,99 ms, katmanlı kurallar en fazla 2,72 ms ve eMPC-HiGHS+safety layer 14,13 ms gerektirmiştir.

\begin{figure}[t]
\centering
\resultfig{figS4_inference_time.pdf}{$N$'ye göre kontrol adımı başına çıkarım süresi.}
\caption{Politika ve güvenlik katmanı dâhil, $N$'ye göre kontrol adımı başına çıkarım süresi.}
\label{fig:inference}
\end{figure}

\restable{T6}{\placeholdertable{tab:t6}{Hesaplama}{T6.tex}}

\subsection{Göreli açıklar, stres düzeyi ve literatür yapılandırması}\label{app:res-sensitivity}

\Tab{tab:ts1}, yumuşak cezalı optimum ve her iki katı optimum için göreli açıkları toplamların oranı olarak raporlamaktadır.
$J^{\sm{pen,LP}}$ çoğu zaman sıfıra yakın ya da negatif olduğundan \method{} yönteminin göreli $J^{\sm{pen}}$ açığı $N=5$ ve $N=10$ için \%100'ü aşmıştır; bu nedenle ana metin mutlak açığı kullanmaktadır.
\Tab{tab:ts2}, $N=20$ ve $N=200$ için $\rho_{\sm{tr}}=3$ kW stres düzeyini kapsamaktadır.
Burada \method{} yönteminin E4 değeri \%55,30 ve \%65,91; Deep Sets+safety layer'ınki ise \%65,14 ve \%73,25 olmuştur.
EV2Gym'in 25 şarj üniteli yapılandırmasında Attention+safety layer without departure bound'un açığı 0,363, katmanlı eMPC-HiGHS'nin 0,456 ve \method{} yönteminin 0,507 olmuştur (\Tab{tab:ts3}).

\restable{TS1}{\placeholdertable{tab:ts1}{Toplamların oranı olarak göreli açıklar}{TS1.tex}}
\restable{TS2}{\placeholdertable{tab:ts2}{$\rho_{\sm{tr}}=3$ kW stres düzeyi}{TS2.tex}}
\restable{TS3}{\placeholdertable{tab:ts3}{EV2Gym'in 25 şarj üniteli literatür yapılandırmasındaki sonuçlar}{TS3.tex}}

\subsection{Birleştirilmiş kısıt ölçütleri}\label{app:res-example}

\Fig{fig:constraintspooled}, her yöntem ve trafo anma düzeyi için test ızgarasının tüm istasyon boyutları üzerinden birleştirilmiş karşılanmayan enerjiyi ve aşırı yük enerjisini göstermektedir.
Aşırı yük enerjisi, \Tab{tab:t3c} içinde katmanlı her yöntemin sıfır aşırı yük adımıyla uyumlu olarak yalnızca güvenlik katmanı içermeyen yöntemlerde görülmektedir.
Konformal varyantlar gösterilmemiştir; bunların eksik enerjisini \Tab{tab:t7} vermektedir.

\begin{figure}[t]
\centering
\resultfig{figS6_constraints_pooled.pdf}{İstasyon boyutları üzerinden birleştirilmiş, yöntem başına kısıt ölçütleri.}
\caption{İstasyon boyutları üzerinden birleştirilmiş, yöntem ve trafo anma düzeyi başına kısıt ölçütleri.}
\label{fig:constraintspooled}
\end{figure}
